\chapter{API HTTP}\label{cap:api}

La interfaz Qadra y cualquier cliente automatizado operan el sistema a través
de la API HTTP que expone el subcomando \texttt{serve} (\cref{cap:cli}). Este
capítulo resume las convenciones comunes a todas las rutas y señala qué
documento del repositorio contiene el contrato completo de cada familia. El
contrato general es \archivo{docs/http-api.md}; el índice por tema,
\archivo{docs/README.md}.

El portal de la autoridad certificadora tiene un contrato propio, servido
por el binario \texttt{qadra-ca} y descrito en \archivo{docs/ca/api.md}
(\cref{cap:ca}).

\section{Base y versionado de rutas}\label{sec:api-base}

Todas las rutas usan el prefijo \texttt{/api/v1}, con una sola excepción:
\texttt{GET /healthz}, que es pública, informa el estado del proceso y queda
fuera del control de admisión. Por omisión el servidor escucha únicamente en
\texttt{127.0.0.1:3000}; la publicación hacia el exterior corresponde al proxy
descrito en el \cref{cap:despliegue}.

Los recursos de un expediente cuelgan de
\texttt{/api/v1/cases/\{case\_id\}}. Las rutas documentales siempre incluyen
el expediente: las antiguas rutas globales \texttt{/api/v1/documents}
responden \texttt{404} y no son alias de las actuales. Los identificadores son
UUID y los cuerpos de petición y respuesta son JSON, salvo la carga de
documentos (binaria o \emph{multipart}) y las descargas de contenido y
evidencia.

\section{Autenticación por sesión con segundo factor}\label{sec:api-auth}

El acceso se obtiene en dos pasos y produce un token opaco, no un JWT:

\begin{enumerate}
  \item \texttt{POST /api/v1/auth/login} verifica correo y contraseña y
        devuelve un desafío MFA con vigencia de 5 minutos.
  \item \texttt{POST /api/v1/auth/mfa/totp}, o bien la ruta
        \archivo{/api/v1/auth/mfa/recovery} con un código de recuperación,
        recibe \texttt{challenge\_token} y \texttt{code}, consume el
        desafío y devuelve \texttt{access\_token}, \texttt{token\_type},
        \texttt{expires\_in\_seconds}, \texttt{user} y los metadatos de
        vigencia de la sesión.
\end{enumerate}

Las peticiones protegidas envían el token en la cabecera
\texttt{Authorization}:

\begin{verbatim}
Authorization: Bearer <access_token>
\end{verbatim}

Con la sesión abierta, \texttt{GET} \archivo{/api/v1/auth/me} devuelve la
identidad vigente, \texttt{GET} \archivo{/api/v1/auth/session} informa los
plazos sin renovarlos y \texttt{POST} \archivo{/api/v1/auth/logout} revoca el
token de inmediato. La primera cuenta de Administrador (rol \texttt{owner})
se crea con \texttt{POST}
\archivo{/api/v1/auth/bootstrap}, ruta pública solo mientras no exista ningún
usuario (\cref{cap:instalacion}).

Conviene tener presentes estas reglas al integrar un cliente o al diagnosticar
un rechazo:

\begin{itemize}
  \item El token consta de 256 bits aleatorios. Redis conserva solo una clave
        derivada con SHA-256, de modo que cerrar sesión lo revoca sin
        almacenar el token en claro.
  \item El rol y la pertenencia a un expediente no viajan en el token: cada
        petición protegida recarga el estado de la cuenta en PostgreSQL y
        confirma la sesión en Redis. Retirar una asignación o desactivar una
        cuenta surte efecto sobre la sesión ya abierta.
  \item La sesión vence a las 24 horas. El vencimiento por inactividad está
        deshabilitado por omisión y se activa con
        \texttt{-{}-session-idle-seconds} o
        \texttt{TT\_SESSION\_IDLE\_SECONDS} (entero de 1 a 86400). Solo
        \texttt{POST /api/v1/auth/activity} amplía ese plazo; las demás
        peticiones no renuevan nada.
  \item Cinco contraseñas rechazadas bloquean el correo durante 15 minutos
        (\texttt{429 account\_locked}); un desafío MFA admite un solo
        intento, aunque sea rechazado.
\end{itemize}

\subsection{Ejemplo de primer paso de acceso}\label{sec:api-ejemplo}

Primer paso con \texttt{curl} y datos ficticios, tomado del contrato:

{\small
\begin{verbatim}
base=http://127.0.0.1:3000
curl --fail-with-body -X POST "$base/api/v1/auth/login" \
  -H 'Content-Type: application/json' \
  --data '{"email":"owner@example.com",
           "password":"correct horse battery staple"}'
\end{verbatim}
}

La respuesta contiene el desafío. El segundo paso envía
\texttt{challenge\_token} y \texttt{code} a la ruta TOTP, y las llamadas
posteriores (por ejemplo \texttt{GET /api/v1/auth/me}) llevan la cabecera
\texttt{Authorization} mostrada arriba. El recorrido completo, incluidos
permisos por rol, evidencia y revocación, se reproduce contra servicios
temporales con \texttt{bash scripts/api-demo.sh}.

\section{Cabeceras}\label{sec:api-cabeceras}

\begin{itemize}
  \item \texttt{Authorization}: esquema \texttt{Bearer}, sin distinguir
        mayúsculas. Se rechazan cabeceras múltiples y tokens con espacios.
  \item \texttt{Content-Type: application/json} en los cuerpos JSON, y
        \texttt{X-Document-Name} con el nombre del archivo en la carga
        binaria de documentos.
  \item \texttt{Cache-Control: no-store} en las respuestas de la API.
  \item \texttt{X-Request-Id} en las respuestas de las rutas MFA: lo genera
        el servidor y permite localizar el intento en los registros
        descritos en \archivo{docs/mfa-access-logs.md}.
  \item \texttt{X-Document-Digest} en la descarga del paquete de evidencia.
\end{itemize}

El actor siempre se toma de la sesión: la API no acepta cabeceras ni campos
que declaren identidad, rol o reloj del cliente. La dirección IP registrada en
auditoría procede del \emph{socket}; \texttt{X-Forwarded-For} solo se
considera con \texttt{TT\_AUDIT\_TRUST\_LOOPBACK\_PROXY=true} y un proxy en
\emph{loopback} (\cref{cap:configuracion}).

\section{Formato de errores}\label{sec:api-errores}

Todos los errores comparten una envoltura estable. El campo \texttt{code} es
el identificador que debe usar un cliente; \texttt{message} es descriptivo:

\begin{verbatim}
{"error":{"code":"permission_denied","message":"permission denied"}}
\end{verbatim}

\begin{table}[H]
  \centering
  \caption{Códigos de estado frecuentes.}\label{tab:api-estados}
  \begin{tabularx}{\textwidth}{L{1.4cm} Y L{5.2cm}}
    \toprule
    \textbf{Estado} & \textbf{Significado} & \textbf{Ejemplos de \texttt{code}} \\
    \midrule
    \texttt{400} & UUID, JSON o parámetro mal formado; campos desconocidos o
                   repetidos. & \texttt{invalid\_json} \\
    \texttt{401} & Credenciales, segundo factor o sesión inválidos.
                 & \texttt{invalid\_session}, \texttt{mfa\_rejected} \\
    \texttt{403} & Rol autenticado sin permiso.
                 & \texttt{permission\_denied} \\
    \texttt{404} & Recurso inexistente, oculto o ajeno al expediente
                   indicado.
                 & \texttt{case\_not\_found}, \texttt{document\_not\_found} \\
    \texttt{409} & Conflicto de revisión, expediente cerrado u operación ya
                   realizada.
                 & \texttt{case\_closed}, \texttt{document\_already\_sealed} \\
    \texttt{413} & Cuerpo mayor que el límite de la ruta.
                 & \texttt{document\_too\_large} \\
    \texttt{422} & Valores sintácticamente correctos pero inválidos.
                 & \texttt{invalid\_input} \\
    \texttt{429} & Ventana de inicio de sesión bloqueada.
                 & \texttt{account\_locked} \\
    \texttt{500} & Fallo interno, sin detalles del \emph{backend}.
                 & \texttt{internal\_error} \\
    \texttt{503} & Presupuesto de peticiones agotado.
                 & \texttt{server\_busy} \\
    \bottomrule
  \end{tabularx}
\end{table}

\section{Límites, paginación y reintentos}\label{sec:api-limites}

\paragraph{Admisión.} El servidor admite por omisión ocho peticiones
simultáneas y dos operaciones bloqueantes; se ajustan con
\texttt{-{}-max-in-flight-requests} y \texttt{-{}-max-blocking-operations}.
Al agotarse el presupuesto, la siguiente petición recibe de inmediato
\texttt{503 server\_busy}. No hay reintentos automáticos ni un tiempo de
espera global.

\paragraph{Tamaño de los cuerpos.} Cada ruta fija su propio límite: 16~KiB en
identidad y alta básica de expedientes, de 8~KiB a 1~MiB en los demás cuerpos
JSON según la familia, y 16~MiB en la carga de documentos. El detalle está en
la sección \emph{Límites HTTP y sobrecarga} de \archivo{docs/http-api.md}.

\paragraph{Paginación.} No existe un esquema único; cada contrato declara el
suyo. Los listados de expedientes y documentos usan \texttt{limit} (de 1 a
100) y \texttt{offset}. Otras colecciones, como el directorio de
participantes, usan un cursor exclusivo (\texttt{after\_id},
\texttt{before\_revision}) y devuelven \texttt{has\_more} con el cursor
siguiente. No ofrecen un total global ni una instantánea entre páginas.

\paragraph{Concurrencia y reintentos.} La API no usa una cabecera de
idempotencia. La repetición segura se resuelve de tres maneras:

\begin{itemize}
  \item Asignar y retirar miembros de un expediente (\texttt{PUT} y
        \texttt{DELETE}) son operaciones idempotentes.
  \item Las ediciones envían la revisión esperada del recurso; una revisión
        desactualizada produce \texttt{409} y el cliente debe revisar los
        datos vigentes antes de reenviar.
  \item Las escrituras con evidencia siguen el patrón \emph{preparar y
        confirmar}: \texttt{POST \ldots/prepare} devuelve el comando
        normalizado y su resumen sin escribir, y la confirmación deja un
        recibo de operación consultable. Ante una respuesta incierta se
        consulta ese recibo; no se repite la escritura a ciegas.
\end{itemize}

\importante{Un \texttt{500} no acredita que la operación se haya revertido:
el registro de acceso se escribe después de la operación de negocio. Concilie
con el recibo o con el estado actual del recurso antes de reenviar.}

\section{Familias de rutas}\label{sec:api-familias}

La \cref{tab:api-familias} relaciona cada familia con el archivo que
contiene sus cuerpos, respuestas, permisos y errores; todos están en el
directorio \archivo{docs/}. Las rutas se abrevian: llevan el prefijo
\texttt{/api/v1}, y \texttt{\{c\}} equivale a
\texttt{/cases/\{case\_id\}}.

{\small
\begin{longtable}{L{3.6cm} L{5.4cm} L{6.1cm}}
  \caption{Familias de rutas y documento que las describe.}
  \label{tab:api-familias}\\
  \toprule
  \textbf{Familia} & \textbf{Base de las rutas} & \textbf{Contrato} \\
  \midrule
  \endfirsthead
  \toprule
  \textbf{Familia} & \textbf{Base de las rutas} & \textbf{Contrato} \\
  \midrule
  \endhead
  \bottomrule
  \endfoot
  Identidad y sesiones
    & \archivo{/auth}, \archivo{/users}
    & \archivo{http-api.md} \\
  Recuperación de contraseña
    & \archivo{/auth/password-reset}
    & \archivo{password-reset-public.md} \\
  Miembros y acceso
    & \archivo{/users}, \archivo{{c}/members}
    & \archivo{members-api.md} \\
  Invitaciones
    & \archivo{/member-invitations}, \archivo{/auth/invitations}
    & \archivo{member-invitations-api.md} \\
  Certificado del Administrador
    & \archivo{/auth/certificate-bindings}, \archivo{/auth/certificate-login}
    & \archivo{http-owner-certificates.md},
      \archivo{owner-certificate-login-http.md} \\
  Expedientes y administración
    & \archivo{/cases}, \archivo{/penal-cases},
      \archivo{/case-administrations}
    & \archivo{http-api.md},
      \archivo{case-administration-api.md} \\
  Etapas procesales
    & \archivo{{c}/stage}
    & \archivo{case-stages-api.md} \\
  Participantes
    & \archivo{{c}/participants}, \archivo{{c}/subjects}
    & \archivo{http-api.md},
      \archivo{typed-participants-api.md} \\
  Documentos y contenido
    & \archivo{{c}/documents}, \archivo{/document-integrity-incidents}
    & \archivo{http-api.md},
      \archivo{document-upload-admission-api.md},
      \archivo{document-content-api.md} \\
  Firma individual
    & \archivo{/identity/signing-certificates},
      \archivo{{c}/documents/.../individual-signing}
    & \archivo{http-individual-signing.md} \\
  Contratos bilaterales
    & \archivo{{c}/contracts}, \archivo{/contract-requests}
    & \archivo{http-contract-signing.md} \\
  Confianza de la evidencia
    & \archivo{/evidence/trust}
    & \archivo{http-api.md} \\
  Audiencias y resultados
    & \archivo{{c}/hearings}, \archivo{/hearings}
    & \archivo{hearings-api.md},
      \archivo{hearing-results-api.md} \\
  Medidas cautelares
    & \archivo{{c}/precautionary-hearings}, \archivo{{c}/measures},
      \archivo{{c}/measure-decisions}
    & \archivo{precautionary-hearings-api.md},
      \archivo{measure-records-api.md},
      \archivo{measure-decisions-api.md} \\
  Resoluciones y notificaciones
    & \archivo{{c}/resolutions}
    & \archivo{procedural-facts-api.md} \\
  Recursos procesales
    & \archivo{{c}/procedural-resources}
    & \archivo{procedural-resources-api.md},
      \archivo{resource-activities-api.md},
      \archivo{resource-deadlines-api.md},
      \archivo{activity-resource-links-api.md} \\
  Calendarios
    & \archivo{/judicial-calendars}
    & \archivo{judicial-calendars-api.md} \\
  Plazos
    & \archivo{{c}/deadlines}, \archivo{/deadline-profiles}
    & \archivo{deadlines-api.md},
      \archivo{deadline-tracking-api.md},
      \archivo{deadline-profiles-api.md} \\
  Agenda
    & \archivo{/agenda}
    & \archivo{agenda-api.md} \\
  Alertas
    & \archivo{/alerts}, \archivo{/alert-preferences}
    & \archivo{alerts-api.md} \\
  Tablero
    & \archivo{/dashboard}
    & \archivo{dashboard-api.md} \\
  Informes
    & \archivo{/case-reports}
    & \archivo{case-reports-api.md} \\
  Auditoría
    & \archivo{/audit/events}, \archivo{/audit/verify}
    & \archivo{audit-events-api.md} \\
\end{longtable}
}

La recuperación pública de contraseña, las invitaciones y el acceso del
Administrador con certificado permanecen deshabilitados mientras no se
configuren de forma expresa; cada contrato indica sus variables. La matriz de permisos por
ruta está en \archivo{docs/access-audit-matrix.md}, y cada petición admitida
deja un evento de acceso en la bitácora de auditoría (\cref{cap:operacion}).
